\subsection{Lean verification}

The proof compiler is Lean 4.32.0, commit
\code{8c9756b28d64dab099da31a4c09229a9e6a2ef35}, matching the repository's
toolchain. The direct bundle theorem, the free-name bound, preservation of free
support through equality interpretation, and the direct seed theorem have been
compiled against the reviewed source. The successful named-bundle and direct
seed reports use only \code{propext}, \code{Classical.choice}, and
\code{Quot.sound}.

The final clean replay compiles 112 project modules serially with no project
object reuse. Every compilation succeeds. The two typed audit clients check
the generic endpoint and both direct/legacy seed APIs, producing eleven
expected axiom reports. The two free-support lemmas use \code{propext} and
\code{Quot.sound}; the other nine audited results use the standard three
axioms. There are no warnings in the changed modules or audit clients; 18
preexisting linter warnings occur in unchanged dependencies. The recorded
elapsed time is 113.52 seconds for this replay, with no baseline speed
comparison implied.

The exact compilation records are described in \code{lean/README.md} and
its accompanying evidence. The interface split is
checked separately from the optional old quotient route, so retaining old
clients cannot silently reintroduce that route into the new K215 import closure.
The source-only census computes the full hypothetical K215 closure of the
overlay, including its unchanged physical certificate modules; it does not
compile those certificates.

\subsection{Rocq component verification}

The local component check uses Coq 8.18.0 with OCaml 4.14.1. Ten selected source
modules are compiled serially, including the shared core, both actual
derivability instances, both compatibility facades, and the unchanged
\code{ZF.L1_LMeta} and \code{K215F.L224Meta} consumers. The audit restates all
142 original theorem signatures. Its 26 printed assumption reports are closed:
12 public propositional rules in each calculus and two meta-theory results.
Recursive \code{coqchk} succeeds for the selected consumers and their imports.

The only compatibility edit in the old meta-theory consumer for this local
check is the import-prefix spelling \code{Stdlib} to \code{Coq}. The packet
does not claim that Coq 8.18 is the production toolchain: the maintained K215
Rocq lane pins Rocq 9.1.0 and stdlib 9.0.0. The supplied component must be
registered in that driver and checked there before adoption. Raw exit codes,
source and object hashes, API logs, and the recursive kernel-check record are
retained in \code{rocq/evidence/}.

\subsection{Independent review and limits of performance claims}

An independent source-and-receipt review verifies the direct argument's
completeness premise, the exact import counts, the moved interface statements,
the 142 Rocq API signatures, ten Rocq source/object records, and the closed
assumption reports. The final Lean replay has its own 112 bound module
records. This review is labeled as a review, not a second compiler
run. Its detailed findings are included with the archive.

The evidence contains per-module elapsed times to reproduce the work's scope.
These are not a controlled before/after build benchmark. No clean-build
speedup or proof-object-size reduction is claimed. Module counts, physical
source lines, source bytes, and number of maintained derivation bodies are
kept as separate metrics.
